\documentclass[11pt]{article}

% Packages
\usepackage[margin=1in]{geometry}
\usepackage{amsmath, amssymb, amsthm}
\usepackage{enumitem}
\usepackage[most]{tcolorbox}
\usepackage{hyperref}
\usepackage{mathtools}

% Theorem environments
\theoremstyle{definition}
\newtheorem{definition}{Definition}[section]
\theoremstyle{plain}
\newtheorem{theorem}[definition]{Theorem}
\theoremstyle{remark}
\newtheorem{remark}[definition]{Remark}

% Dark axiom box style
\newtcolorbox{axiombox}{
  colback=black!85!white,
  coltext=white,
  colframe=black!90!white,
  boxrule=0pt,
  arc=4pt,
  left=10pt, right=10pt, top=8pt, bottom=8pt,
}

% Intuition box style
\newtcolorbox{intuitionbox}[1][]{
  colback=blue!5!white,
  colframe=blue!40!black,
  boxrule=0.8pt,
  arc=4pt,
  title={\textbf{Intuition}},
  left=8pt, right=8pt, top=6pt, bottom=6pt,
  #1
}

% Operators
\DeclareMathOperator{\rank}{rank}

\begin{document}

% Title
\begin{center}
  {\Huge \textbf{Symmetric Matrix Eigenvalue Properties}} \\[12pt]
  {\Large EECS 127} \\[24pt]
\end{center}

\vspace{10pt}

% Problem Statement
\begin{axiombox}
\textbf{Problem.} \; Let $A \in \mathbb{R}^{n \times n}$ be a real symmetric matrix, i.e.\ $A = A^T$.

\begin{enumerate}[label=(\alph*), itemsep=6pt]
  \item Prove that every eigenvalue of $A$ is real.
  \item Prove that eigenvectors corresponding to distinct eigenvalues are orthogonal.
\end{enumerate}
\end{axiombox}

\vspace{20pt}

%% ---- Part (a) ---- %%
\section*{Part (a): Eigenvalues are real}

\emph{Idea: compute $v^*Av$ two ways, using $A^* = A$.}

\begin{proof}
Let $\lambda \in \mathbb{C}$, $v \in \mathbb{C}^n \setminus \{0\}$, $Av = \lambda v$.

Left-multiply by $v^*$:
\[
  v^* (Av) = v^* (\lambda v) = \lambda \, (v^* v).
\]

On one hand, take the conjugate transpose of $v^* A v$. Since $A$ is real, $\bar{A} = A$, so $A^* = A^T = A$:
\[
  (v^* A v)^* = v^* A^* v = v^* A v,
\]
so $v^* A v = \overline{v^* A v}$, i.e.\ $v^* A v \in \mathbb{R}$.

On the other hand, conjugate $\lambda\,(v^*v)$. Note $v^* v = \|v\|^2 > 0$ is real, so:
\[
  \overline{\lambda \, (v^* v)} = \bar\lambda \, \overline{(v^* v)} = \bar\lambda \, (v^* v).
\]

Combining: $v^* A v$ is real, so $v^* A v = \overline{v^* A v}$. But $v^* A v = \lambda\,(v^*v)$, so:
\[
  \lambda \, (v^* v) = \overline{\lambda \, (v^* v)} = \bar\lambda \, (v^* v).
\]

Since $v \neq 0$, we have $v^* v > 0$, so divide both sides:
\[
  \boxed{\lambda = \bar\lambda,}
\]
hence $\lambda \in \mathbb{R}$. \qedhere
\end{proof}

\vspace{20pt}

%% ---- Part (b) ---- %%
\section*{Part (b): Eigenvectors for distinct eigenvalues are orthogonal}

\emph{Idea: compute $u_1^T A\, u_2$ two ways, using $A = A^T$.}

\begin{proof}
Let $\lambda_1, \lambda_2$ be distinct eigenvalues with eigenvectors $u_1, u_2 \in \mathbb{R}^n$:
\[
  A u_1 = \lambda_1 u_1, \qquad A u_2 = \lambda_2 u_2, \qquad \lambda_1 \neq \lambda_2.
\]

Way A --- act on the right with $A u_2 = \lambda_2 u_2$:
\[
  u_1^T (A u_2) = u_1^T (\lambda_2 u_2) = \lambda_2 \, (u_1^T u_2).
\]

Way B --- act on the left. Since $A = A^T$, rewrite $u_1^T A$:
\[
  u_1^T A = u_1^T A^T = (A u_1)^T = (\lambda_1 u_1)^T = \lambda_1\, u_1^T.
\]
So:
\[
  u_1^T (A u_2) = \lambda_1\, u_1^T u_2 = \lambda_1 \, (u_1^T u_2).
\]

Equate Way A and Way B:
\[
  \lambda_2 \, (u_1^T u_2) = \lambda_1 \, (u_1^T u_2).
\]

Rearrange:
\[
  (\lambda_1 - \lambda_2)\, u_1^T u_2 = 0.
\]

Since $\lambda_1 \neq \lambda_2$, we have $(\lambda_1 - \lambda_2) \neq 0$, so divide:
\[
  \boxed{u_1^T u_2 = 0.} \qedhere
\]
\end{proof}

\end{document}
